\lecture{6.7}{Evaluating the Super Learner}
\label{lec:6-7}

\subsection{\texttt{CV.SuperLearner}}

The function \texttt{SuperLearner()} uses internal cross-validation to choose ensemble
weights. That internal risk is useful for comparing library algorithms, but it is not by
itself an external evaluation of the final Super Learner procedure. To estimate the
out-of-sample risk of the Super Learner itself, use \texttt{CV.SuperLearner()}.

\begin{implementationbox}
\begin{lstlisting}[language=R]
set.seed(1234)

cv_sl <- CV.SuperLearner(
  Y = Y,
  X = data.frame(W1 = W1, W2 = W2, A = A),
  family = gaussian(),
  method = "method.CC_LS",
  SL.library = c("SL.glm", "SL.mean", "SL.quadglm"),
  cvControl = list(V = 2),
  innerCvControl = list(list(V = 2), list(V = 2))
)
\end{lstlisting}
\end{implementationbox}

\subsection{Nested Cross-Validation}

Evaluation of the Super Learner requires nested cross-validation. There are two
levels:
\begin{itemize}
\item outer folds estimate the out-of-sample risk of the whole Super Learner
procedure;
\item inner folds, inside each outer training set, choose the Super Learner ensemble
weights.
\end{itemize}

\begin{figure}[H]
\centering
\begin{tikzpicture}[
  node distance=1.1cm,
  box/.style={draw, rounded corners=2pt, minimum width=3.0cm, minimum height=.75cm},
  small/.style={draw, rounded corners=2pt, minimum width=2.35cm, minimum height=.62cm},
  arr/.style={-{Latex[length=2.5mm]}, thick}
]
\node[box] (data) {Full data};
\node[box, below left=1cm and .8cm of data] (outertrain) {Outer training};
\node[box, below right=1cm and .8cm of data] (outervalid) {Outer validation};
\node[small, below=.9cm of outertrain] (inner) {Inner CV\\chooses weights};
\node[small, below=.9cm of outervalid] (risk) {Evaluate\\risk};
\draw[arr] (data) -- (outertrain);
\draw[arr] (data) -- (outervalid);
\draw[arr] (outertrain) -- (inner);
\draw[arr] (inner) -- node[above] {trained SL} (risk);
\draw[arr] (outervalid) -- (risk);
\end{tikzpicture}
\caption{Nested cross-validation evaluates the entire Super Learner procedure.}
\label{fig:ch6-nested-cv}
\end{figure}

If the same folds used to choose weights were also used to evaluate the final Super
Learner, the risk estimate would be optimistic. Nested cross-validation avoids this by
holding out outer validation observations from every part of fitting and weight
selection.

\subsection{Cross-Validated Risk Estimates}

\texttt{CV.SuperLearner()} returns risk estimates for:
\begin{itemize}
\item each candidate algorithm;
\item the discrete Super Learner;
\item the ensemble Super Learner.
\end{itemize}

For continuous outcomes with squared-error loss, the risk is estimated mean squared
prediction error:
\[
\widehat R_{CV}
=
\frac{1}{n}\sum_{i=1}^n
\{Y_i-\hat Q_{-v(i)}(A_i,W_i)\}^2,
\]
where $\hat Q_{-v(i)}$ denotes the prediction rule trained without the outer validation
fold containing observation $i$.

\subsection{Confidence Intervals for Risk}

Plots from \texttt{CV.SuperLearner()} often display estimated risk and approximate
confidence intervals. These intervals quantify uncertainty in the estimated predictive
risk, not uncertainty in a causal effect. They are useful for judging whether apparent
risk differences are meaningful or mostly noise.

If two algorithms have very similar risk estimates with overlapping intervals, it is not
reasonable to claim that one is clearly better. The ensemble can still use both, or put
weight on only one, depending on the metalearner optimization.

\subsection{Discrete Super Learner}

The discrete Super Learner is the single candidate with the best cross-validated risk.
When evaluating results, it is helpful to compare:
\[
\text{best single candidate},\qquad
\text{discrete Super Learner},\qquad
\text{ensemble Super Learner}.
\]
In some datasets, the ensemble improves on all individual candidates. In others, the
best single candidate performs about as well as the ensemble.

\subsection{Comparing Super Learner to Candidate Algorithms}

The following questions guide interpretation:
\begin{itemize}
\item Does the ensemble have lower estimated risk than the best single candidate?
\item Are risk differences large relative to their uncertainty?
\item Which algorithms receive nonzero weight?
\item Are any algorithms unstable, slow, or producing warnings?
\item Does the selected library behave differently for $\Qbar$ than for $g$?
\end{itemize}

\begin{keyidea}
\texttt{SuperLearner()} fits an ensemble for use in analysis. \texttt{CV.SuperLearner()}
evaluates how well the entire ensemble-building procedure predicts new data.
\end{keyidea}

\begin{estimationbox}
For causal effect estimation, nested evaluation is a diagnostic for nuisance-function
prediction. It does not prove that the resulting causal estimator is unbiased, because
causal validity still depends on identification assumptions and adequate positivity.
\end{estimationbox}
